\chapter{Functions of Two Variables}\label{ch:b1:multivar}

The year ends with a first walk into higher dimension: functions
$f(x, y)$ of two real variables. Everything generalizes --- limits,
continuity, derivatives, extrema --- but each notion gains a twist:
limits can be approached along every direction at once, derivatives
split into partials, and the gradient points the way uphill. The
full theory (differentials, general $\R^n$, submanifolds) belongs to
the second year; here we set the vocabulary and the first honest
theorems.

\section{The plane $\R^2$; continuity}

\begin{definition}\label{def:b1:multivar:topology}
On $\R^2$, use the Euclidean norm $\norm{(x,y)} = \sqrt{x^2 + y^2}$
(\cref{ch:b1:euclid}). Open balls, neighborhoods, \emph{open
subsets} of $\R^2$ are defined exactly as in
\cref{ch:b1:topology}, with balls in place of intervals. A function
$f \colon U \to \R$ ($U \subseteq \R^2$ open) is \emph{continuous
at} $a \in U$ when
\[
\forall\varepsilon > 0,\ \exists\delta > 0, \quad
\norm{X - a} \leq \delta \implies \abs{f(X) - f(a)} \leq
\varepsilon,
\]
with the same sequential characterization as in one variable.
Sums, products, quotients and compositions with continuous
one-variable functions preserve continuity; the coordinate maps are
continuous, hence so are polynomials in $(x,y)$.
\end{definition}

\begin{example}[The radial trap]\label{ex:b1:multivar:trap}
Let $f(x, y) = \dfrac{xy}{x^2 + y^2}$ for $(x,y) \neq (0,0)$, $f(0,
0) = 0$. Along each axis, $f = 0 \to 0$; but along the diagonal $y
= x$, $f(x, x) = \frac12 \not\to 0$. \emph{No limit at the origin}:
approaching along every line, and even finding the same limit along
each, is not enough (here the line limits disagree; worse examples
agree along all lines yet fail along a parabola,
\cref{exo:b1:multivar:3}). Continuity in each variable separately
does not imply continuity.
\end{example}

\section{Partial derivatives}

\begin{definition}\label{def:b1:multivar:partial}
The \emph{partial derivatives}\index{partial derivative} of $f$ at
$(a, b)$ are the one-variable derivatives along the axes:
\[
\frac{\partial f}{\partial x}(a,b)
= \lim_{h \to 0} \frac{f(a + h, b) - f(a,b)}{h},
\qquad
\frac{\partial f}{\partial y}(a,b)
= \lim_{k \to 0} \frac{f(a, b + k) - f(a,b)}{k}.
\]
$f$ is of class $C^1$ on $U$ when both exist and are continuous on
$U$. The \emph{gradient}\index{gradient} is $\nabla f(a,b) =
\bigl(\frac{\partial f}{\partial x},\, \frac{\partial f}{\partial
y}\bigr)(a,b)$.
\end{definition}

\begin{theorem}[$C^1$ implies a tangent plane]\label{thm:b1:multivar:c1}
Let $f$ be $C^1$ on $U$ and $(a,b) \in U$. Then, as $(h, k) \to
(0,0)$:
\[
f(a + h, b + k) = f(a, b)
+ h\,\frac{\partial f}{\partial x}(a,b)
+ k\,\frac{\partial f}{\partial y}(a,b)
+ o\bigl(\norm{(h,k)}\bigr) .
\]
In particular $f$ is continuous, and the graph $z = f(x,y)$ has at
each point the tangent plane read off the formula.
\end{theorem}

\begin{proof}
Move one coordinate at a time:
\[
f(a+h, b+k) - f(a,b)
= \bigl[f(a+h, b+k) - f(a, b+k)\bigr] + \bigl[f(a, b+k) -
f(a,b)\bigr].
\]
By the one-variable mean value theorem
(\cref{thm:b1:derivative:mvt}), the first bracket is $h\,
\frac{\partial f}{\partial x}(a + \theta h,\, b + k)$ for some
$\theta \in \intoo{0}{1}$, and the second is $k\,\frac{\partial
f}{\partial y}(a,\, b + \theta' k)$. Continuity of the partials at
$(a,b)$ lets us write each as (value at $(a,b)$) $+$
(error $\to 0$); the total error is $h\,\varepsilon_1 +
k\,\varepsilon_2 = o(\norm{(h,k)})$ since $\abs h, \abs k \leq
\norm{(h,k)}$.
\end{proof}

\begin{theorem}[Chain rule]\label{thm:b1:multivar:chain}
Let $f$ be $C^1$ on $U$ and $t \mapsto (x(t), y(t))$ be $C^1$ from
an interval into $U$. Then $g(t) = f\bigl(x(t), y(t)\bigr)$ is
$C^1$, with\index{chain rule!two variables}
\[
g'(t) = x'(t)\,\frac{\partial f}{\partial x}\bigl(x(t),y(t)\bigr)
+ y'(t)\,\frac{\partial f}{\partial y}\bigl(x(t),y(t)\bigr)
= \bigl\langle \nabla f,\ (x', y')\bigr\rangle .
\]
\end{theorem}

\begin{proof}
Apply \cref{thm:b1:multivar:c1} at $(x(t), y(t))$ with $(h, k) =
(x(t+s) - x(t),\, y(t+s) - y(t))$: as $s \to 0$, $h = s\,x'(t) +
o(s)$, $k = s\,y'(t) + o(s)$, $\norm{(h,k)} = O(s)$; divide by $s$.
Continuity of $g'$ follows from that of all ingredients.
\end{proof}

\begin{remark}[Reading the gradient]
Along a unit direction $u$, the chain rule on $t \mapsto f(a + tu)$
gives the \emph{directional derivative} $\langle \nabla f(a),
u\rangle$: maximal when $u$ points along $\nabla f(a)$ (by
Cauchy--Schwarz, \cref{thm:b1:euclid:cs}). The gradient is the
direction of steepest ascent, and it is orthogonal to the level
curves $\{f = c\}$ (differentiate $f$ along a curve drawn in a
level set: the chain rule gives $\langle\nabla f,\,
\text{tangent}\rangle = 0$).
\end{remark}

\begin{theorem}[Schwarz]\label{thm:b1:multivar:schwarz}
If $f$ is of class $C^2$ (partials of partials exist and are
continuous), then\index{Schwarz's theorem}
\[
\frac{\partial^2 f}{\partial x\,\partial y}
= \frac{\partial^2 f}{\partial y\,\partial x} .
\]
\end{theorem}
\admitted

\section{Local extrema}

\begin{theorem}[Critical points]\label{thm:b1:multivar:critical}
If $f$ ($C^1$ on the open set $U$) has a local extremum at $(a,b)
\in U$, then $\nabla f(a,b) = (0,0)$: the point is
\emph{critical}\index{critical point}.
\end{theorem}

\begin{proof}
The one-variable functions $x \mapsto f(x, b)$ and $y \mapsto f(a,
y)$ have interior local extrema at $a$, resp.\ $b$:
\cref{prop:b1:derivative:fermat} kills both partials.
\end{proof}

\begin{method}[Second-order test (Monge notation)]\label{met:b1:multivar:monge}
At a critical point of a $C^2$ function, set
\[
r = \frac{\partial^2 f}{\partial x^2},
\qquad
s = \frac{\partial^2 f}{\partial x \partial y},
\qquad
t = \frac{\partial^2 f}{\partial y^2}
\qquad (\text{values at the point}).
\]
\begin{itemize}
  \item If $rt - s^2 > 0$: local extremum --- minimum for $r > 0$,
        maximum for $r < 0$;
  \item if $rt - s^2 < 0$: no extremum (a \emph{saddle
        point}\index{saddle point});
  \item if $rt - s^2 = 0$: the test is silent; examine directly.
\end{itemize}
(Justification --- a Taylor--Young expansion at order $2$ and the
sign study of the quadratic form $r h^2 + 2shk + tk^2$ --- is
carried out in the second year; the test is used here as a working
tool.)
\end{method}

\begin{example}\label{ex:b1:multivar:extrema}
$f(x,y) = x^3 + y^3 - 3xy$. Critical points: $\nabla f = (3x^2 -
3y,\; 3y^2 - 3x) = 0$ gives $y = x^2$ and $x = y^2$, so $x = x^4$:
$x \in \{0, 1\}$: points $(0,0)$ and $(1,1)$.

Second derivatives: $r = 6x$, $s = -3$, $t = 6y$. At $(0,0)$: $rt -
s^2 = -9 < 0$: saddle. At $(1,1)$: $rt - s^2 = 36 - 9 > 0$, $r = 6
> 0$: local minimum, $f(1,1) = -1$. (Not global: $f(x, 0) = x^3 \to
-\infty$.)
\end{example}

\begin{omfigure}
\begin{tikzpicture}
\begin{axis}[omaxis, width=7.6cm, height=6.6cm,
    xmin=-1.6, xmax=2.1, ymin=-1.6, ymax=2.1,
    xtick={-1,1}, ytick={-1,1},
    xlabel={$x$}, ylabel={$y$}]
  % level curves of f = x^3+y^3-3xy: f = -1 (min), f=-0.5, 0, 1
  \addplot[omDef, thick, domain=-1.35:1.9, samples=80]
    {x^2} node[pos=0.94, left, font=\small]
    {$\frac{\partial f}{\partial x} = 0$};
  \addplot[omThm, thick, domain=-1.35:1.9, samples=80]
    ({x^2}, {x}) node[pos=0.97, below right, font=\small]
    {$\frac{\partial f}{\partial y} = 0$};
  \fill (axis cs:0,0) circle (1.8pt)
    node[below left, font=\small] {saddle};
  \fill (axis cs:1,1) circle (1.8pt)
    node[below right, font=\small] {min};
\end{axis}
\end{tikzpicture}

{\small The two curves $\nabla f = 0$ for $f = x^3 + y^3 - 3xy$
intersect at the critical points $(0,0)$ (saddle) and $(1,1)$
(local minimum).}
\end{omfigure}

\section{Exercises}

\begin{exercise}[$\star$]\label{exo:b1:multivar:1}
Compute the partial derivatives:
$f(x,y) = x^2 y + \eu^{xy}$; $\;g(x,y) = \ln(x^2 + y^2)$ (on
$\R^2\setminus\{0\}$); $\;h(x,y) = \arctan\frac yx$ (on $x > 0$).
\end{exercise}

\begin{exercise}[$\star$]\label{exo:b1:multivar:2}
Study the continuity at $(0,0)$ (with value $0$ there) of:
\[
f(x,y) = \frac{x^2 y}{x^2 + y^2},
\qquad
g(x,y) = \frac{xy}{x^2 + y^2},
\qquad
h(x,y) = \frac{x^3 + y^3}{x^2 + y^2}.
\]
\emph{(Polar coordinates $x = \rho\cos\theta$, $y = \rho\sin\theta$
help: bound by a function of $\rho$ alone when possible.)}
\end{exercise}

\begin{exercise}[$\star\star$]\label{exo:b1:multivar:3}
Let $f(x,y) = \dfrac{x^2 y}{x^4 + y^2}$ ($f(0,0) = 0$). Prove that
$f$ has limit $0$ at the origin \emph{along every straight line},
but that $f\bigl(x, x^2\bigr) = \frac12$: $f$ is not continuous at
$(0,0)$.
\end{exercise}

\begin{exercise}[$\star$]\label{exo:b1:multivar:4}
Verify Schwarz's theorem by hand on $f(x, y) = x^3 y^2 + \sin(xy)$.
\end{exercise}

\begin{exercise}[$\star\star$]\label{exo:b1:multivar:5}
Let $f$ be $C^1$ on $\R^2$ and $g(t) = f(\cos t, \sin t)$. Express
$g'(t)$ via the chain rule. Deduce that $f$ restricted to the unit
circle attains extrema at points where $\nabla f$ is parallel to
the radius vector.
\end{exercise}

\begin{exercise}[$\star\star$]\label{exo:b1:multivar:6}
Find and classify the critical points of:
\[
f(x, y) = x^2 + xy + y^2 - 3x,
\qquad
g(x, y) = x^2 - y^2 + 4y,
\qquad
h(x, y) = x^4 + y^4 - 2(x - y)^2 .
\]
\emph{(For $h$, the determinant test is silent at the origin:
examine $h(x, x)$ and $h(x, -x)$.)}
\end{exercise}

\begin{exercise}[$\star\star$]\label{exo:b1:multivar:7}
A function $f$ is \emph{harmonic} when $\frac{\partial^2 f}{\partial
x^2} + \frac{\partial^2 f}{\partial y^2} = 0$. Check that $x^2 -
y^2$, $xy$, $\eu^x\cos y$ and $\ln(x^2 + y^2)$ (off the origin) are
harmonic.
\end{exercise}

\begin{exercise}[$\star\star\star$]\label{exo:b1:multivar:8}
Find the point of the plane $\{x + 2y - z = 4\} \subseteq \R^3$
closest to the origin, two ways: by orthogonal projection
(\cref{ch:b1:euclid}), and by minimizing the two-variable function
$f(x, y) = x^2 + y^2 + (x + 2y - 4)^2$ obtained by eliminating $z$.
\end{exercise}

\begin{exercise}[$\star\star\star$]\label{exo:b1:multivar:9}
(Laplacian in polar coordinates, first contact) Let $f$ be $C^2$ on
$\R^2 \setminus \{0\}$ and \emph{radial}: $f(x, y) =
\varphi\bigl(\sqrt{x^2+y^2}\bigr)$ with $\varphi$ of class $C^2$ on
$\intoo{0}{+\infty}$. Prove that
\[
\frac{\partial^2 f}{\partial x^2} + \frac{\partial^2 f}{\partial
y^2} = \varphi''(\rho) + \frac{\varphi'(\rho)}{\rho},
\qquad \rho = \sqrt{x^2 + y^2},
\]
and find all radial harmonic functions on the punctured plane.
\end{exercise}
